\clearpage
\thispagestyle{empty}
\begin{center}
    \bfseries\fontsize{14}{16}\selectfont ABSTRACT
\end{center}
\vspace{16pt}

\noindent
The \textit{Capacitated Vehicle Routing Problem} (CVRP) is a classical combinatorial
optimization problem in which a minimum-cost set of routes must be designed to serve
customers with known demands by means of a homogeneous fleet of capacity-constrained
vehicles. This report addresses the CVRP exclusively from the standpoint of
exact mathematical programming methods. The state of the art of \textit{Branch and X}
algorithms for the problem is reviewed, covering \textit{Branch and Bound},
\textit{Branch and Cut}, \textit{Branch and Price} and \textit{Branch, Cut and Price},
and the adoption of the classical scheme by Laporte and Nobert (1983) is justified as
the methodological basis of the work. The document presents the formal definition of
the CVRP, its mathematical formulation as an integer linear program over an undirected
graph, and the resolution procedure based on the iterative generation of
capacity inequalities over an initial relaxation with degree constraints. The computational implementation, developed in Python using PuLP and the
COIN-OR CBC solver, was evaluated on five standard CVRPLIB instances (Sets E
and A). The method certified optimality in four of them (0\% gap) and, on the
fifth, reached a feasible solution within 0.91\% of the optimum under the time
budget, empirically exposing the scalability limit of the classical scheme.

\vspace{10pt}
\noindent\textbf{Keywords:} CVRP; combinatorial optimization; integer linear
programming; exact methods; \textit{Branch and Bound}; CVRPLIB.